\scanpage{205}
est nul, où \deleted{\unsure{nous}}\inserted{$\{a\}$} est le schéma ouvert
induit sur $a$, donc $\{a\}$ est dominé
par $\,{}_{(mn)}\Gamma$. Cela prouve bien que
la suite des $\,{}_n\Gamma$ est loc. const, que son
sup est un ouvert, $\Gamma^\tau$ qui est caractérisé
comme on l'a dit. Que cet ouvert
soit \deleted{invariant}\inserted{un sous-groupe} est immédiat, et résulte aussi du

\deleted{N.B.} \uline{Corollaire}. Pour tout $S$ \deleted{\unsure{local}}\inserted{qu. cpct} sur $k$,
$\Gamma^\tau(S)={}$ torsion de $\Gamma(S)$. \deleted{\unsure{Et les fibres|Et la plus}}

\deleted{\uline{Corollaire}. $\Gamma^\tau$ est le plus petit sous-groupe
fermé de $\Gamma$ tel que $\Gamma/\Gamma^\tau$ soit séparable sans \unsure{tor.}}

Si $G$ est un \deleted{\unsure{sché}} $k$-groupe
\deleted{quelconq} loc de type fini, on
désigne par $G^\tau$ l'image inverse de
$\Gamma^\tau$ par $G\unsure{-|\to}G/G_0=\Gamma$.

\uline{Proposition} $G^{\deleted{\unsure{0|\tau}}}_0$ \inserted{(resp. $G_{\unsure{\tau|\varepsilon}}^{\deleted{\unsure{\tau|0}}}$)} est le plus petit sous-groupe
fermé de $G$ tel que $G/G_0$ soit
\uline{discret} \uline{séparable} [resp. discret \deleted{séparable}\inserted{sans}
torsion]. C'est aussi le sup. \deleted{\unsure{$\bigcup_n U_n$}}\inserted{\unsure{(croissant)|croissant)} des sous-groupes [ouverts; \unsure{ou : ouvert} \unsure{``\quad''} }]
de $G$ \unsure{(de} \deleted{\unsure{type}} \unsure{finis,} \unsure{resp. les}
% slanted
\margin{à prouver \ldots}
$G$ est dit \uline{``sans torsion''} si
$(G/G_0)^\tau=\{e\}$, \deleted{\unsure{Alors}} $G^\tau$ est dit le
\uline{sous-groupe de torsion de $G$}
% archive sheet number: 204
